\documentclass[11pt,twocolumn,a4paper]{article}

% ─── Page layout ─────────────────────────────────────────────────────
\usepackage[margin=2cm, columnsep=0.6cm]{geometry}

% ─── Essential packages ──────────────────────────────────────────────
\usepackage[utf8]{inputenc}
\usepackage[T1]{fontenc}
\usepackage{lmodern}
\usepackage{amsmath,amssymb,amsthm,mathtools}
\usepackage{bm}
\usepackage{graphicx}
\usepackage{booktabs}
\usepackage{array}
\usepackage{multirow}
\usepackage{enumitem}
\usepackage{float}
\usepackage{siunitx}
\usepackage{microtype}
\usepackage{caption}
\usepackage{natbib}
\usepackage{xcolor}
\usepackage[colorlinks=true,linkcolor=blue!60!black,%
           citecolor=blue!60!black,urlcolor=blue!60!black]{hyperref}

\captionsetup{font=small,labelfont=bf,textfont=it}

% ─── Theorem environments ────────────────────────────────────────────
\newtheorem{theorem}{Theorem}[section]
\newtheorem{lemma}[theorem]{Lemma}
\newtheorem{proposition}[theorem]{Proposition}
\newtheorem{corollary}[theorem]{Corollary}
\theoremstyle{definition}
\newtheorem{definition}[theorem]{Definition}
\theoremstyle{remark}
\newtheorem{remark}[theorem]{Remark}
\newtheorem{example}[theorem]{Example}

% ─── Custom commands ─────────────────────────────────────────────────
\newcommand{\Rens}{R_{\mathrm{ens}}}
\newcommand{\Kc}{K_{c}}
\newcommand{\sigw}{\sigma_{\omega}}
\newcommand{\Hcal}{\mathcal{H}}
\newcommand{\Ccal}{\mathcal{C}}
\newcommand{\Tcal}{\mathcal{T}}
\newcommand{\Mcal}{\mathcal{M}}
\newcommand{\Gcal}{\mathcal{G}}
\newcommand{\Fcal}{\mathcal{F}}
\newcommand{\Ocal}{\mathcal{O}}
\newcommand{\Pcal}{\mathcal{P}}
\newcommand{\Bcal}{\mathcal{B}}
\newcommand{\Acal}{\mathcal{A}}
\newcommand{\dP}{\Delta P}
\newcommand{\Tref}{T_{\mathrm{ref}}}
\newcommand{\trec}{t_{\mathrm{rec}}}
\newcommand{\fref}{f_{\mathrm{ref}}}
\newcommand{\indep}{\perp\!\!\!\perp}

\title{\textbf{Football as a Partially Observable Bioreactor:\\[2pt]
\large Backward Biomechanical Accounting via the Ball Observation Operator}}

\author{Kundai Farai Sachikonye\\[2pt]
\small AIMe Registry for Artificial Intelligence\\
\small \texttt{kundai.sachikonye@bitspark.com}}
\date{\today}

\begin{document}

\twocolumn[%
\maketitle%
\begin{abstract}
\noindent
Quantitative analysis of association football has historically split between
single-player kinematic studies, which discard team-level coupling, and
event-level statistics, which discard biomechanics. We propose a third
framework: the team is modelled as a partially observable bioreactor of
eleven coupled biomechanical cells, and the ball is the primary observation
operator that samples one cell at a time at high fidelity while broadcast
imagery serves as a low-resolution ambient indicator of the remaining ten.
The cellular ensemble is given a coupled-oscillator dynamics on the
Kuramoto manifold; the team-level Kuramoto order parameter
$\Rens \in {[0,1]}$ supplies a continuous coherence measure with an
analytically sharp phase-locked regime above which a single \emph{theoretical
player} ---a federation-wide cell registry---suffices to describe the
collective. Duels between opposing players are formalised as cognate
oscillator pairs whose coupling is computed jointly; this is the
irreducible biomechanical accounting unit. A measurement substrate is
proposed in which only the lower body is tracked: pitch-perimeter
advertising panels function as planar calibration targets whose homography
solves both metric ground-plane recovery and player/spectator
segmentation in a single step. Goal sequences are reconstructed
\emph{backwards} as morphism compositions in a category whose objects are
player-roles and whose morphisms are passes; the generators of this
category constitute a team's tactical-biomechanical signature. We argue
that the FIFA World Cup, with 104 matches and 48 distinct collective
templates, supplies a near-ideal corpus for the framework. The output is
a biomechanical \emph{ledger}, not a predictive engine.
\end{abstract}
\vspace{0.5em}]

% ═══════════════════════════════════════════════════════════════════
\section{Introduction}
% ═══════════════════════════════════════════════════════════════════

\subsection{The measurement problem}

Two extant traditions dominate quantitative analysis of association
football. The first is single-player kinematics: stride frequency,
ground-reaction force, joint kinematics, and metabolic estimates
\citep{cavagna1977,kram1990,novacheck1998,winter2009}, typically gathered
from instrumented running on treadmills or from short field bouts. This
tradition produces precise biomechanics but discards the coupling between
players that defines team sport. The second tradition is event-level
statistics: expected goals, expected possessions, pass networks
\citep{fernandez2018,decroos2019}, typically aggregated from
broadcast-derived event streams. This tradition captures team-level
patterns but discards biomechanics entirely. A third strand---automated
broadcast tracking systems \citep{theagarajan2021,dorazio2009}---attempts
to bridge the gap but is bounded by the spatial resolution of broadcast
footage, in which players typically occupy 30--80 pixels in height and
robust full-skeleton pose extraction is unreliable.

The present work proposes a complementary framework. Rather than treating
a team as eleven independent measurement targets, we treat it as a single
partially observable system with eleven coupled internal cells, sampled by
a privileged probe---the ball---that occupies one cell at a time. The
framework is biomechanical (the cells are oscillator-coupled bodies, not
abstract tokens) and team-level (the ball's trajectory through cells is
the principal observable). Our central claim is that this framing
extracts more biomechanically meaningful information from
already-available broadcast footage than either tradition does alone, and
does so with a measurement substrate that side-steps the pose-resolution
problem.

\subsection{Bioreactor analogy}

A continuous-flow bioreactor contains a parallel population of cells whose
states evolve simultaneously and continuously, but which an experimenter
typically samples one at a time at high resolution while monitoring the
bulk vessel through low-resolution ambient indicators (turbidity,
dissolved oxygen, pH) \citep{danckwerts1953,levenspiel1999,bailey1986}.
The reactor's state at time $t$ is the joint distribution over all cell
states; the experimenter's belief about that state is the conditional
distribution given the sample history and the ambient observations. This
is a classical example of partial observability
\citep{kalman1960,luenberger1971,kaelbling1998}.

We claim football has precisely this structure. The eleven players of a
team are eleven cells; the ball is the high-resolution probe (only one
player possesses it at any instant, and the biomechanics of that player
during possession is densely measurable from foot-contact timing); the
broadcast camera is the low-resolution ambient indicator (the off-ball
players are visible enough to estimate formation and pressing density,
but not at the per-foot-contact fidelity afforded by the ball-bearer).
The analogy is more than rhetorical: in Sections \ref{sec:obsop} and
\ref{sec:backward} we show that the bioreactor sampling formalism
\citep{danckwerts1953} supplies the correct mathematical machinery for
the backward reconstruction of team trajectories from goals.

\subsection{Contributions}

This paper makes the following contributions:

\begin{enumerate}[leftmargin=1.4em]
  \item A formal observation-operator model of a football team as a
    partially observable cellular system, with the ball as primary
    probe (Section~\ref{sec:obsop}).
  \item A coupled-oscillator dynamics for the team in which the
    Kuramoto order parameter $\Rens$ supplies a continuous, regime-banded
    measure of tactical coherence (Section~\ref{sec:kuramoto}).
  \item A categorical formalisation of the \emph{theoretical player}
    abstraction, identified with a federation-wide cell registry
    (Section~\ref{sec:theoretical}).
  \item A formalisation of duels as cognate oscillator pairs, supplying
    the irreducible biomechanical accounting unit
    (Section~\ref{sec:cognate}).
  \item A measurement substrate using lower-body tracking against
    advertising-panel homography (Section~\ref{sec:substrate}).
  \item A backward goal-reconstruction algorithm that recovers a
    team's morphism category from its scoring trajectories
    (Section~\ref{sec:backward}).
  \item An argument that the FIFA World Cup is a near-ideal dataset
    for the framework, with concrete computable quantities specified
    (Section~\ref{sec:wc}).
\end{enumerate}

The procedure is biomechanical accounting---a structured tally of what
bodies actually did, organised by team coherence regime, duel
participation, and morphism composition. It is not a predictive engine
and not a tactics recommender.

% ═══════════════════════════════════════════════════════════════════
\section{The Observation-Operator Framework}
\label{sec:obsop}
% ═══════════════════════════════════════════════════════════════════

\subsection{Cellular ensemble}

\begin{definition}[Team-bioreactor]
\label{def:bioreactor}
A \emph{team-bioreactor} at instant $t$ is the tuple
$\Bcal(t) = (\{X_i(t)\}_{i=1}^{N}, \Phi(t), \Psi(t))$ where:
\begin{itemize}[leftmargin=1.4em]
  \item $\{X_i(t)\}_{i=1}^{N}$ is the collection of \emph{biomechanical
    cell states}, one per active player, with $N = 11$ in a regular
    association-football team.
  \item $\Phi(t) \in \mathbb{R}^{2 \times N}$ is the \emph{formation
    matrix} of player pitch positions.
  \item $\Psi(t)$ is the bulk pressing/spacing state of the unsampled
    cells, extracted from the ambient observation.
\end{itemize}
\end{definition}

The cell state $X_i$ is the biomechanical configuration of player $i$:
its gait phase, joint configuration, instantaneous ground-reaction force,
and any other locally measurable quantity. The configuration space of an
individual cell is denoted $\Hcal_i$ and is assumed Polish; the joint
configuration space is $\Hcal = \prod_{i=1}^{N} \Hcal_i$.

\subsection{The ball as observation operator}

\begin{definition}[Possession indicator]
\label{def:possession}
The \emph{possession indicator} is a piecewise-constant function
$\pi : [0,T] \to \{1, \ldots, N, *\}$ where $\pi(t) = i$ if player $i$
is in contact with the ball at time $t$, and $\pi(t) = *$ otherwise
(the ball is in flight or otherwise unpossessed).
\end{definition}

\begin{definition}[Ball observation operator]
\label{def:obs}
The \emph{ball observation operator} is the partial map
$\Ocal_{\mathrm{ball}} : \Hcal \times \{1,\ldots,N,*\} \to \Hcal_i \cup \{\emptyset\}$
defined by
\begin{equation}
  \Ocal_{\mathrm{ball}}\!\left(\{X_i\},\pi\right) =
  \begin{cases}
    X_{\pi} & \pi \in \{1,\ldots,N\}, \\
    \emptyset & \pi = *.
  \end{cases}
\end{equation}
\end{definition}

The operator $\Ocal_{\mathrm{ball}}$ projects the full bioreactor state
to the single cell in possession of the ball. The remaining ten cells
continue to evolve under their own intrinsic dynamics during possession,
but the observation does not directly record them; their states are
recovered, if at all, from the ambient bulk indicator $\Psi$.

\begin{proposition}[Sample-trajectory completeness]
\label{prop:complete}
Over a finite time horizon $[0,T]$ the ball observation operator
$\Ocal_{\mathrm{ball}}$ produces a piecewise-defined trajectory in the
disjoint union $\bigsqcup_i \Hcal_i$ whose total measure is
$\lambda(\{t : \pi(t) \neq *\})$ and whose pieces are indexed by
$\pi(t)$.
\end{proposition}

\begin{proof}
The set $\{t \in [0,T] : \pi(t) \neq *\}$ is a finite disjoint union
of intervals $\{I_k\}_{k=1}^{K}$ on which $\pi$ is constant. On each
$I_k$ the operator records $X_{\pi(I_k)}(t)$ for $t \in I_k$, which is
a continuous trajectory in $\Hcal_{\pi(I_k)}$. The disjoint union of
these pieces, indexed by the constant value of $\pi$ on $I_k$, gives
the result. The total measure follows by additivity.
\end{proof}

\subsection{Bulk ambient observation}

\begin{definition}[Ambient observation]
\label{def:ambient}
The \emph{ambient observation} is a map $\Acal :
\{\text{video frame}\} \to \mathbb{R}^{q}$ extracting a finite-dimensional
summary of the bulk team state (formation centroid, pressing density,
team width, depth, compactness) at low spatial resolution. The
ambient observation is generally available at frame rate and at all
times $t \in [0,T]$, irrespective of $\pi(t)$.
\end{definition}

The pair $(\Ocal_{\mathrm{ball}}, \Acal)$ defines a partial observability
structure in the standard sense \citep{kaelbling1998}. The state estimator
maintains a posterior over the unsampled cells given the history of both
channels; this is a standard Bayesian filtering problem with the wrinkle
that the high-fidelity channel switches between cells (the ball changes
possession) over time.

\begin{figure*}[!htbp]
    \centering
    \includegraphics[width=\textwidth]{figures/panel1_observation_operator.png}
    \caption{\textbf{Observation-operator sampling on the pitch-bioreactor.}
    (\textbf{A})~Three-dimensional rendering of a schematic football pitch
    with eleven player positions in a 4--3--3 formation (cyan) and the
    ball trajectory (amber) lifted into a vertical ``sampled cell''
    coordinate. The ball touches one player at a time at high fidelity,
    rising into the sampled-cell coordinate at each contact and travelling
    between cells in arcs; the remaining ten cells evolve at $z=0$ in the
    bioreactor's bulk and are not directly sampled.
    (\textbf{B})~Possession indicator $\pi(t)$ over a 80-second window.
    Each cyan block is a possession bout of a single player; the grey
    gaps mark $\pi(t)=*$ (ball in flight or unpossessed). The function
    is piecewise constant, with bout boundaries triggering re-targeting
    of the high-fidelity sample channel.
    (\textbf{C})~Distribution of possession-bout durations across 400
    simulated bouts. The mean bout duration is $1.43$~s, well above the
    Nyquist threshold $1/(2B)=0.20$~s for a biomechanical signal of
    bandwidth $B=2.5$~Hz (red dashed line), confirming that typical
    elite-football bouts Nyquist-sample gait biomechanics for any player
    whose touches are not micro-deflections.
    (\textbf{D})~Per-player total possession time over a single match,
    coloured by the Nyquist threshold (green: above 8~s, red: below).
    The accounting framework requires sufficient cumulative sample mass
    to reconstruct each player's biomechanical trajectory; players with
    insufficient touches (defenders rarely on the ball) rely more
    heavily on the ambient observation channel.}
    \label{fig:obsop}
\end{figure*}

\subsection{Sampling theorem analogue}

A residence-time-style argument bounds the information available from
the ball channel alone.

\begin{theorem}[Cellular sampling bound]
\label{thm:sampling}
Let $\overline{\tau}_i$ denote the expected possession duration of
player $i$ within a match of total length $T$. The total ball-channel
sample time allocated to player $i$ satisfies
\begin{equation}
  \mathbb{E}\!\left[\lambda(\pi^{-1}\{i\})\right] = \overline{\tau}_i \cdot n_i,
\end{equation}
where $n_i$ is the expected number of possession bouts for player $i$.
If the per-player intrinsic biomechanical bandwidth is $B_i$, the
ball-channel suffices to reconstruct $X_i$ up to Nyquist
\citep{shannon1948} only if $\overline{\tau}_i \geq 1/(2B_i)$ on
average across bouts.
\end{theorem}

\begin{proof}
Each possession bout is a finite-duration sample of $X_i$. The classical
sampling theorem \citep{shannon1948,cover2006} requires a sample interval
of at most $1/(2B_i)$ to reconstruct a band-limited signal of bandwidth
$B_i$. The expected total sample mass is the product of bout count and
mean bout duration by linearity of expectation.
\end{proof}

For typical football biomechanics $B_i$ is dominated by stride frequency
($\sim$1.5--3 Hz), so $\overline{\tau}_i \geq 0.17$--$0.33$~s suffices.
Mean possession bouts in elite football exceed 1~s, so the ball channel
alone Nyquist-samples gait biomechanics for any player whose touches are
not micro-deflections. Sub-Nyquist players (defenders rarely touching
the ball) require greater reliance on the ambient channel.

% ═══════════════════════════════════════════════════════════════════
\section{Coupled-Oscillator Team Dynamics}
\label{sec:kuramoto}
% ═══════════════════════════════════════════════════════════════════

\subsection{Players as phase oscillators}

The locomotion of a running player is fundamentally oscillatory: the
gait cycle is a stereotyped periodic motion with stride frequency
$f_i$ in the 1--3~Hz range under match conditions
\citep{novacheck1998,cavagna1977}. We assign to each player a phase
variable $\phi_i \in [0, 2\pi)$ advancing at the player's
instantaneous angular gait frequency $\omega_i = 2\pi f_i$.

\begin{definition}[Team Kuramoto model]
\label{def:kuramoto}
The collective dynamics of the eleven phase oscillators is governed
by the Kuramoto equation \citep{kuramoto1984}
\begin{equation}
  \dot{\phi}_i = \omega_i + \frac{K}{N}\sum_{j=1}^{N}
                 \sin(\phi_j - \phi_i),
  \label{eq:kuramoto}
\end{equation}
with natural frequencies $\omega_i$ drawn from a symmetric unimodal
distribution $g(\omega)$ of standard deviation $\sigw$, and coupling
strength $K \geq 0$ representing the strength of tactical
co-ordination between teammates.
\end{definition}

The coupling $K$ is not a physical force in the conventional sense; it
is the effective strength with which a teammate's stride pattern,
posture, or anticipated trajectory modulates an individual player's
gait. Coupling arises through visual feedback, learned response
patterns, and tactical drilling, and it is the modelling proxy for what
match analysts call ``cohesion'' or ``shape.''

\subsection{Order parameter and regimes}

\begin{definition}[Team order parameter]
\label{def:order}
The \emph{team order parameter} is the complex average
\begin{equation}
  \Rens(t)\,e^{i\psi(t)} = \frac{1}{N}\sum_{j=1}^{N} e^{i\phi_j(t)},
  \label{eq:order}
\end{equation}
with $\Rens(t) \in [0,1]$ and mean phase $\psi(t) \in [0, 2\pi)$.
$\Rens = 0$ indicates a uniform random phase distribution across the
team; $\Rens = 1$ indicates perfect phase alignment.
\end{definition}

The classical Kuramoto bifurcation
\citep{kuramoto1984,strogatz2000,acebron2005} establishes a critical
coupling
\begin{equation}
  \Kc = \frac{2}{\pi g(0)}
  \label{eq:Kc}
\end{equation}
above which a partially phase-locked solution
$\Rens^{*} = \sqrt{1 - \Kc/K}$ emerges via a supercritical pitchfork.
For symmetric unimodal distributions of width $\sigw$, $\Kc$ scales
linearly with $\sigw$; teams whose players have a narrow distribution
of natural stride frequencies require less tactical coupling to achieve
phase coherence than teams with wide frequency spread.

\begin{definition}[Coherence regimes]
\label{def:regimes}
The interval $[0,1]$ is partitioned into five operationally meaningful
regimes by the values $\Rens \in \{0.30, 0.50, 0.80, 0.95\}$:
\begin{enumerate}[leftmargin=1.4em]
  \item \emph{Turbulent} ($\Rens < 0.30$): players act as independent
    oscillators; no team-level coherence.
  \item \emph{Aperture-dominated} ($0.30 \leq \Rens < 0.50$): partial
    correlations exist but a unified template is unwarranted.
  \item \emph{Cascade} ($0.50 \leq \Rens < 0.80$): a dominant cluster
    exists, but multiple sub-units operate semi-independently.
  \item \emph{Coherent} ($0.80 \leq \Rens < 0.95$): a single dominant
    cluster carries the team; minor stragglers remain.
  \item \emph{Phase-locked} ($\Rens \geq 0.95$): full coherence; the
    team acts as a single template (Section~\ref{sec:theoretical}).
\end{enumerate}
\end{definition}

These regime boundaries are derived from the mean-field self-consistency
relation by inverting $\Rens^* = \sqrt{1-\Kc/K}$ at the values
$\Rens \in \{0.3, 0.5, 0.8, 0.95\}$, giving $K/\Kc$ ratios of
$1.10, 1.33, 2.78, 10.25$ respectively. The boundary between regimes is
a continuous transition in the underlying coupling potential
\citep{strogatz2000} that nevertheless produces operationally sharp
changes in the kind of collective behaviour the team can sustain.


\begin{figure*}[!htbp]
    \centering
    \includegraphics[width=\textwidth]{figures/panel2_bifurcation.png}
    \caption{\textbf{Kuramoto bifurcation and the five coherence regimes.}
    (\textbf{A})~Three-dimensional surface of the team order parameter
    $R_{\mathrm{ens}}(K, t)$ across a sweep of coupling strengths
    $K/K_{c}\in[0.5,8]$ and time $t\in[0,20]$~s, for an ensemble of
    $N=40$ Cauchy-distributed oscillators. The surface rises from the
    incoherent floor at low $K$ to a fully synchronised plateau at high
    $K$, with the bifurcation at $K=K_{c}=2\gamma$ visible as the steep
    ridge.
    (\textbf{B})~Equilibrium order parameter $R_{\mathrm{ens}}$ versus
    $K/K_{c}$ for $N=50$ simulated ensembles (cyan dots) against the
    mean-field prediction $R^{*}=\sqrt{1-K_{c}/K}$ (violet line). The
    bifurcation at $K=K_{c}$ (red dashed) separates the incoherent
    state ($R\approx 0$) from the partially phase-locked state
    ($R>0$); the simulated curve matches the mean-field prediction
    above threshold within finite-$N$ tolerance.
    (\textbf{C})~Sinusoidal phase trajectories $\sin(\phi_i(t))$ for
    eleven oscillators at sub-critical ($K=0.7\,K_{c}$, grey) and
    super-critical ($K=4\,K_{c}$, violet, offset) coupling. The
    sub-critical trajectories drift independently; the super-critical
    trajectories synchronise to a common envelope, illustrating the
    team-template emergence at coherent coupling.
    (\textbf{D})~Coherence regime classification. Each of the five
    operational regimes (turbulent, aperture, cascade, coherent,
    phase-lock) appears as a horizontal jitter scatter of $200$
    sample points at its centre $R$ value with the regime mean marked
    by the black bar. Dashed grey lines at $R\in\{0.30, 0.50, 0.80,
    0.95\}$ delineate the regime boundaries.}
    \label{fig:bifurcation}
\end{figure*}

\subsection{Estimation from observable data}

\begin{proposition}[$\Rens$ from foot-contact timing]
\label{prop:rest}
Let $\hat{\phi}_i(t)$ be an estimate of player $i$'s gait phase computed
by linear interpolation between successive observed foot-contact events
on player $i$. Then
\begin{equation}
  \hat{\Rens}(t) = \left| \frac{1}{N}\sum_{j=1}^{N} e^{i\hat{\phi}_j(t)} \right|
  \label{eq:Rens-hat}
\end{equation}
is a consistent estimator of $\Rens(t)$ as the foot-contact sampling
rate increases.
\end{proposition}

\begin{proof}
For each player, linear interpolation in phase between foot-contact
events reconstructs $\phi_i(t)$ exactly when the gait is purely periodic
and degrades smoothly with cycle-to-cycle variability \citep{winter2009}.
The plug-in estimator \eqref{eq:Rens-hat} is the empirical mean of unit
vectors over the team and converges to the true team mean
\eqref{eq:order} by the law of large numbers applied to the
finite-population mean.
\end{proof}

\subsection{Phase-lock and the validity of templating}

The phase-locked regime is operationally privileged because it is the
regime in which speaking of ``one collective body'' becomes precise.

\begin{theorem}[Templating threshold]
\label{thm:template}
Suppose the team is phase-locked with $\Rens \geq 0.95$. Then for any
two players $i, j$ and any time horizon $\Delta t$ short relative to
the slow dynamics of $K$, the cell states $X_i, X_j$ differ from the
team-template state $X^{*}$ only by an offset bounded in $\Hcal$ by
$\epsilon(\Rens) = 2\sin\!\left((1-\Rens)^{1/2}/\sqrt{2}\right)$.
\end{theorem}

\begin{proof}
At $\Rens = 1$ all phases coincide and the cells share a state by
definition. For $\Rens < 1$, the mean-field result \citep{strogatz2000}
gives the per-oscillator phase deviation
$|\phi_i - \psi| \leq \arccos(\Rens)$. The deviation in $\Hcal$
between two phase-locked cells is bounded by
$\|X_i - X^{*}\| \leq L\,|\phi_i - \psi|$ for $L$ the local
Lipschitz constant of the cell embedding; combining with the standard
$\arccos$ bound and the small-angle approximation
$\arccos(\Rens) \leq \pi(1-\Rens)$ for $\Rens \in [1/2,1]$ yields the
stated bound up to a constant absorbed into the global Lipschitz scale.
\end{proof}

The theorem makes precise the slogan ``the team is one body when it is
phase-locked.'' Above the threshold, a single template suffices to
describe the team's biomechanics to within a controlled error; below it,
separate per-player accounting is required.

% ═══════════════════════════════════════════════════════════════════
\section{The Theoretical Player}
\label{sec:theoretical}
% ═══════════════════════════════════════════════════════════════════

\subsection{Categorical formulation}

The slogan ``the team is one player'' is formalised most cleanly in
categorical terms \citep{maclane1971,awodey2010,spivak2014}.

\begin{definition}[Role category]
\label{def:rolecat}
The \emph{role category} of a team is a small category $\Mcal$ whose
\emph{objects} are player-roles (e.g.\ left-back, deep-lying playmaker,
target striker, transition-six) and whose \emph{morphisms}
$f: r \to r'$ are admissible pass-transitions from role $r$ to role
$r'$, with composition given by sequential passing and the identity
$\mathrm{id}_r$ given by retention of possession.
\end{definition}

The role category is not the literal squad: two distinct human players
occupying the same role contribute morphisms to the same object. The
category encodes the team's tactical identity at a level of abstraction
above personnel.

\begin{definition}[Theoretical player]
\label{def:theoretical}
The \emph{theoretical player} of a team is the role category $\Mcal$
together with a \emph{cell registry} $\Ccal : \mathrm{Mor}(\Mcal) \to
\mathrm{Cells}$ assigning to each morphism a biomechanical cell
(stride pattern, plant timing, release kinematics, etc.) characteristic
of executing that morphism.
\end{definition}

The pair $(\Mcal, \Ccal)$ is the team's biomechanical-tactical signature.
Two distinct teams whose theoretical players are isomorphic as
categories (and whose cell registries agree on the isomorphism) play the
``same football'' even with different personnel. The Spanish national
team and FC Barcelona under a particular sporting director, for example,
plausibly share a theoretical player up to category isomorphism.

\begin{figure*}[!htbp]
    \centering
    \includegraphics[width=\textwidth]{figures/panel3_templating.png}
    \caption{\textbf{Templating threshold: the team is one template only
    when the phase distribution is concentrated.}
    (\textbf{A})~Three-dimensional rosette: 50 oscillator phases plotted
    on a unit circle (cos\,$\phi$, sin\,$\phi$) at five values of
    $R_{\mathrm{ens}}$ stacked vertically. At low $R_{\mathrm{ens}}$
    (bottom) the phases are uniformly distributed around the circle;
    at high $R_{\mathrm{ens}}$ (top) they concentrate at a single mean
    phase. The geometric concentration is the operational meaning of
    ``one collective body.''
    (\textbf{B})~Empirical RMS phase deviation versus $1-R_{\mathrm{ens}}$
    on log--log axes. The theoretical bound $\sqrt{2(1-R_{\mathrm{ens}})}$
    (violet line) and the simulated values (cyan dots) coincide to
    machine precision across four decades of $1-R_{\mathrm{ens}}$,
    verifying the rigorous RMS form of the templating threshold theorem.
    (\textbf{C})~Density-normalised phase distributions at three
    $R_{\mathrm{ens}}$ levels (0.50 amber, 0.85 cyan, 0.99 violet).
    The distributions narrow continuously as $R_{\mathrm{ens}}$
    approaches unity, demonstrating that the team-template
    approximation becomes increasingly tight in the phase-locked regime.
    (\textbf{D})~Empirical RMS phase deviation (cyan) and theoretical
    bound $\sqrt{2(1-R_{\mathrm{ens}})}$ (violet) across the coupling
    range $K/K_{c}\in[1, 32]$. Both quantities span more than a decade
    and track each other across the full range, confirming that the
    bound is tight throughout the coherent and phase-locked regimes.}
    \label{fig:templating}
\end{figure*}

\subsection{Goal sequences as morphism compositions}

\begin{definition}[Goal trajectory]
\label{def:goaltraj}
A \emph{goal trajectory} is a finite composition
\begin{equation}
  G = f_n \circ f_{n-1} \circ \cdots \circ f_1 : r_0 \to r_n
  \label{eq:gtraj}
\end{equation}
in $\Mcal$ ending at the distinguished terminal object $r_n =
\mathrm{goal}$, where each $f_k$ is a pass-morphism executed by the team
in real time.
\end{definition}

Working backwards from a scored goal is then the operation of factoring
$G$ in $\Mcal$, recovering the morphism sequence $(f_1, \ldots, f_n)$.
This factorisation is what football analysts intuitively perform when
they ``trace back'' a goal.

\begin{proposition}[Backward unique up to template]
\label{prop:backward}
If the team is phase-locked, the morphism factorisation of $G$ is
unique modulo automorphisms of the role category, and uniquely
recovers a sub-sequence of the cell registry $\Ccal$.
\end{proposition}

\begin{proof}
At phase-lock, by Theorem~\ref{thm:template}, all eleven players share
a state up to bounded offset, so the cell registry $\Ccal$ is shared
across players and the morphism factorisation depends only on the
sequence of roles traversed. Distinct factorisations would correspond
to distinct morphism paths from $r_0$ to $r_n$; under coherent execution
only one path is materially realised in time. Automorphisms of $\Mcal$
that permute roles indistinguishably (e.g.\ symmetric left/right
duplicates of a position) absorb the remaining ambiguity.
\end{proof}

\subsection{Generators}

A team's morphism category is typically generated by a much smaller set
of irreducible play patterns.

\begin{definition}[Tactical generators]
\label{def:gen}
The \emph{tactical generators} of $\Mcal$ are a minimal subset
$\{g_\alpha\} \subset \mathrm{Mor}(\Mcal)$ such that every morphism in
$\Mcal$ is a composition of generators.
\end{definition}

\begin{proposition}[Finite generation in practice]
\label{prop:fingen}
For an elite club or international side, $\Mcal$ is typically generated
by between 8 and 30 distinct play patterns identifiable across multiple
matches.
\end{proposition}

The claim is empirical and is consistent with the observation that
match analysts routinely identify a small finite set of ``set pieces,''
``build-up patterns,'' and ``transition motifs'' as the recognisable
basis of a team's play \citep{hughes2005,memmert2017}. The framework
makes this implicit basis explicit and quantifiable: a team's
generators are the irreducible morphisms whose compositions span all
observed goal trajectories.

% ═══════════════════════════════════════════════════════════════════
\section{Duels as Cognate Oscillator Pairs}
\label{sec:cognate}
% ═══════════════════════════════════════════════════════════════════

\subsection{Coupled-pair subsystem}

Football is fundamentally adversarial: at every instant of contested
play, at least one attacker is paired with at least one defender. The
irreducible biomechanical unit is therefore not a single player but a
pair.

\begin{definition}[Duel]
\label{def:duel}
A \emph{duel} between attacker $a$ and defender $b$ is the
two-oscillator subsystem
\begin{align}
  \dot{\phi}_a &= \omega_a + K_{\mathrm{ab}} \sin(\phi_b - \phi_a) + \xi_a(t), \\
  \dot{\phi}_b &= \omega_b - K_{\mathrm{ab}} \sin(\phi_b - \phi_a) + \xi_b(t),
  \label{eq:duel}
\end{align}
active while attacker and defender are within physical contestation
range. $K_{\mathrm{ab}}$ is the instantaneous duel coupling, large
during contact and zero outside the contestation window.
$\xi_a, \xi_b$ are local stochastic perturbations.
\end{definition}

The opposite sign of the coupling on the attacker and defender reflects
the adversarial structure: attacker seeks to evade, defender seeks to
match, so the pair is anti-synchronously coupled on average.

\begin{proposition}[Cognate-pair coupling]
\label{prop:cognate}
The duel oscillators are \emph{cognate}: their states are co-determined
during the contestation window, and neither is meaningfully accountable
in isolation. Booking the pair as a single biomechanical event respects
the genuine coupling of ground-reaction force and segmental torque
across the contact \citep{novacheck1998,winter2009}.
\end{proposition}

\subsection{Degenerate duels}

A duel may degenerate: the defender's coupling half is null or distant.

\begin{definition}[Degenerate duel]
\label{def:degen}
A duel is \emph{degenerate} if $K_{\mathrm{ab}} = 0$ at the moment of
attacker action, i.e.\ the defender's reach is insufficient to couple
the pair. The corresponding accounting entry is a single-cell action
with no paired partner.
\end{definition}

The degenerate case is the unmarked attacker: a high-value state for
the attacking team, signalled by the absence of a defender half, not by
the presence of any positive measurement on the defender side. The
absence is the signal.

\subsection{Charge-coupled load accounting}

Ground-reaction force and torque transferred through contact during a
duel propagate along the connected skeletal-muscular network of both
bodies. A naive sum of independent per-body biomechanical work would
double-count the contact load.

\begin{theorem}[Pair conservation of contact load]
\label{thm:loadcons}
Let $W_a, W_b$ denote the mechanical work done by attacker and
defender in a contested action, and $W_{\mathrm{c}}$ the work
transferred through the inter-body contact. Then
\begin{equation}
  W_{\mathrm{total}}^{\mathrm{pair}} = W_a + W_b - W_{\mathrm{c}},
  \label{eq:Wpair}
\end{equation}
not the sum $W_a + W_b$.
\end{theorem}

\begin{proof}
By Newton's third law applied at the contact, the impulse delivered by
one body to the other is equal and opposite. The work expended by one
body that becomes kinetic or elastic energy of the other is reckoned
twice if both per-body works are summed independently. Subtracting the
shared contact work $W_{\mathrm{c}}$ yields the genuine total
mechanical work expended by the pair as a unit.
\end{proof}

The theorem motivates pair-level accounting: the duel is the unit at
which contact-coupled load is correctly summed.

\begin{figure*}[!htbp]
    \centering
    \includegraphics[width=\textwidth]{figures/panel4_duels.png}
    \caption{\textbf{Duels as cognate oscillator pairs and the pair
    conservation of contact load.}
    (\textbf{A})~Three-dimensional phase trajectory of a single duel:
    attacker position $x_{a}$, defender position $x_{b}$, and contact
    force $F(t)$ over a $0.30$~s contestation window. Coloured points
    are sampled along the trajectory by time (plasma colourmap); the
    grey shadow at $F=0$ projects the spatial dynamics. The trajectory
    rises from a no-contact state through a peak contact force and
    relaxes back as the bodies separate.
    (\textbf{B})~Velocities of attacker ($v_{a}$, cyan) and defender
    ($v_{b}$, amber) through the duel, with the contact force
    superimposed on a secondary axis (red). The attacker decelerates
    and the defender accelerates symmetrically during the contact peak,
    illustrating the co-determination of the two cognate halves: neither
    velocity profile is meaningful without the other.
    (\textbf{C})~Pair conservation of contact load. The five bars
    show $|\Delta KE_{a}|$ (cyan), $|\Delta KE_{b}|$ (amber), the
    contact-transferred work $W_{c}$ (red), the naive per-body sum
    $W_{a}+W_{b}$ (grey, double-counts the contact transfer), and the
    pair sum $W_{a}+W_{b}-W_{c}$ (violet). The naive sum exceeds the
    correct pair sum by exactly $W_{c}$, demonstrating that booking
    duels at the pair level (Theorem~4.3) avoids the
    contact-double-counting error of independent per-player accounting.
    (\textbf{D})~Scatter of naive sum versus pair sum across $30$
    randomised duels (varied initial velocities, masses, contact
    stiffness, and metabolic inputs). Points fall along a line
    parallel to but below the identity diagonal, with the
    perpendicular offset equal to $W_{c}$ (colour-encoded); the
    mean gap is $1421\pm 496$~J. The systematic offset is the
    framework's prediction.}
    \label{fig:duels}
\end{figure*}

% ═══════════════════════════════════════════════════════════════════
\section{Measurement Substrate}
\label{sec:substrate}
% ═══════════════════════════════════════════════════════════════════

\subsection{Lower-body tracking principle}

Broadcast football footage presents two practical difficulties for
biomechanical extraction: (i) players occupy only 30--80 pixels in
height in wide-angle frames, degrading skeletal pose estimators
\citep{cao2017,bazarevsky2020,newell2016}; (ii) the spectator stand
contains thousands of seated humans whose upper bodies are
indistinguishable from players' at low resolution. Both problems are
addressed by restricting tracking to the lower body.

\begin{proposition}[Lower-body discriminability]
\label{prop:lower}
Spectators in elite-stadium broadcast contexts are uniformly
(i) seated, with lower bodies obscured by stadium seating or by
front-row barriers, and (ii) positioned behind the pitch-perimeter
advertising panels that occlude any visible portion of seated
lower bodies. Consequently, any visible lower body in the broadcast
frame, in front of the perimeter panel line, can be classified as a
player by construction.
\end{proposition}

\subsection{Advertising panels as planar calibration targets}

Modern football stadia universally install LED advertising panels along
the pitch perimeter. These panels are planar, of standard height
($\approx 1$~m), of known geometric dimensions, and carry
high-contrast content suitable for feature-based homography estimation
\citep{hartley2003,zhang2000}.

\begin{theorem}[Panel-homography duality]
\label{thm:panelhom}
Solving the planar homography $H$ that rectifies the visible
advertising-panel content to its known canonical layout
simultaneously:
\begin{enumerate}[leftmargin=1.4em]
  \item Recovers the pitch ground-plane to image-plane projection (the
    panel base coincides with the touch-line by stadium construction),
    yielding metric pitch coordinates for any point on the visible
    ground.
  \item Provides the player/spectator boundary: any visible lower body
    on the pitch side of the rectified panel line is a player; any
    upper body or torso visible above the panel is a spectator,
    occluded below by the panel itself.
\end{enumerate}
\end{theorem}

\begin{proof}
A planar homography in three-dimensional projective space has eight
degrees of freedom and is uniquely determined by four corresponding
points between the image and the canonical layout
\citep{hartley2003,zhang2000}. The advertising panel content provides
substantially more than four trackable features per frame; the
homography is therefore over-determined and robustly estimable
\citep{fischler1981}. The panel base sits at known stadium-specific
metric coordinates along the touch-line; the homography that rectifies
the panel rectifies the ground plane by construction. The
above/below classification follows from the panel's physical
occlusion of any seated lower body and from the spatial constraint
that players are below the panel top.
\end{proof}

\begin{remark}
LED panel content rotates among advertisers every several seconds, but
this does not degrade the homography estimate. Indeed, more graphic
variety supplies more trackable features per frame; homography is
solved on whatever content is currently displayed.
\end{remark}

\subsection{Foot-contact extraction}

The privileged biomechanical primitive in the lower-body substrate is
the foot-contact event: the instant at which a foot strikes the pitch
ground plane.

\begin{definition}[Foot-contact event]
\label{def:fce}
A \emph{foot-contact event} for player $i$ is the pair
$e_{i,k} = (t_{i,k}, p_{i,k})$ where $t_{i,k}$ is the time of the
$k$-th detected ground contact and $p_{i,k} \in \mathbb{R}^2$ is the
pitch metric coordinate of the foot, computed from the panel
homography.
\end{definition}

Detection of foot-contact requires only that the moving lower body
intersect the calibrated ground plane in pixel coordinates and
exhibit the characteristic stationary-then-moving pattern of a foot
strike \citep{novacheck1998}; full skeletal pose is unnecessary. The
foot-contact event stream $\{e_{i,k}\}$ is the principal datum of the
framework.

\subsection{Transcript-driven identity disambiguation}

Match commentary supplies a sparse, low-latency stream of named events
(``Bellingham finds Saka on the right''). This transcript is not
frame-accurate but it disambiguates identity at coarse temporal
resolution. Given metric foot-contact positions and a commentary
event naming the ball-bearer, identity assignment reduces to a
nearest-foot lookup at the event time, which is robust against
typical commentary lag of 1--3 seconds.

% ═══════════════════════════════════════════════════════════════════
\section{Backward Reconstruction from Goals}
\label{sec:backward}
% ═══════════════════════════════════════════════════════════════════

\subsection{The inverse problem}

A scored goal is a terminal event in the morphism category $\Mcal$.
The data available immediately upstream of the goal consist of:
(i) the ball-channel sample trajectory
$X_{\pi(t)}(t)$ for the few seconds leading up to the goal;
(ii) the ambient observation $\Acal(t)$ over the same interval;
(iii) the transcript identification of the named players involved.

The inverse problem is to recover the morphism factorisation
\eqref{eq:gtraj}, the corresponding cell registry entries, and the
off-ball cell states $\{X_j(t)\}_{j \neq \pi(t)}$ during the build-up.

\subsection{Reconstruction algorithm}

\begin{enumerate}[leftmargin=1.4em]
  \item \textbf{Possession segmentation.} Partition the lead-up
    interval into maximal sub-intervals on which $\pi(t)$ is constant.
    Each sub-interval is a possession bout of one named player.
  \item \textbf{Sample-trajectory extraction.} On each bout, extract
    the foot-contact event stream for the ball-bearer; classify
    bouts by their cell-trajectory in the player's local biomechanical
    cell partition.
  \item \textbf{Morphism identification.} Identify each pass (transition
    between bouts) as a morphism in $\Mcal$ by matching the
    sending-bout's cell at release with the receiving-bout's cell at
    first contact.
  \item \textbf{Off-ball cell inference.} For each unsampled player at
    each instant, infer their cell state from the ambient
    observation $\Acal$ (formation, run direction, marking pair)
    constrained by the team's theoretical-player registry.
  \item \textbf{Generator recognition.} Match the recovered morphism
    sequence against the team's known tactical generators
    (Definition~\ref{def:gen}); the goal is decomposed into a
    composition of generators.
\end{enumerate}

\subsection{Aggregate ledger}

\begin{definition}[Goal ledger entry]
\label{def:ledger}
A \emph{goal ledger entry} is the tuple
$L_G = (G, (f_1,\ldots,f_n), \{c_k\}, \{X_j^{(\mathrm{off})}\}, \Rens^{(G)})$
comprising the morphism composition, the per-step cell registry
entries, the inferred off-ball states, and the team coherence
$\Rens$ over the build-up.
\end{definition}

Aggregating $L_G$ over all goals scored by a team across a tournament
yields the empirical generator set and the empirical
phase-coherence distribution conditional on scoring. Both are
biomechanical-tactical signatures of the team that are inaccessible
to either single-player biomechanics or event-level statistics alone.

\begin{figure*}[!htbp]
    \centering
    \includegraphics[width=\textwidth]{figures/panel5_sampling_recovery.png}
    \caption{\textbf{Cellular sampling and backward morphism
    identifiability.}
    (\textbf{A})~Three-dimensional surface of relative reconstruction
    MSE over the joint parameter space (duty fraction, signal
    bandwidth) for a player sampled by the ball channel at cadence
    $3$~Hz. The red curve traces the Nyquist boundary $f_{s}=2B$
    (equivalently $\mathrm{duty}=2B/\mathrm{cadence}$) projected onto
    the $z=0$ plane. The surface is high in the sub-Nyquist corner
    (low duty, high bandwidth) and falls toward zero across the
    super-Nyquist region.
    (\textbf{B})~Relative MSE versus effective sample rate $f_{s}$
    (in Hz, log-scale on $y$). The reconstruction error drops by more
    than two orders of magnitude as $f_{s}$ crosses the Nyquist rate
    $2B=1.0$~Hz (red dashed). Sub-Nyquist sampling produces MSE near
    or above unity; super-Nyquist sampling produces MSE near zero.
    (\textbf{C})~Backward morphism recovery rate versus team coherence
    $R_{\mathrm{ens}}$ during build-up. The recovery curve rises
    monotonically from $\sim$35\% at $R_{\mathrm{ens}}=0.10$ (turbulent)
    to $\sim$98\% at $R_{\mathrm{ens}}=0.95$ (phase-locked). The
    vertical lines mark the boundaries of the coherent regime
    (grey dotted, $R=0.80$) and phase-lock (red dashed, $R=0.95$);
    above coherent, the morphism factorisation is identifiable to
    above $85\%$ accuracy.
    (\textbf{D})~Generator confusion matrix at $R_{\mathrm{ens}}=0.9$.
    Rows are true generator labels and columns are the recovered
    label; cell values are conditional probabilities. The diagonal
    dominance ($>0.92$ across all four generators) confirms that
    in the coherent regime the framework correctly identifies which
    of the team's tactical generators executed a given goal build-up,
    with off-diagonal confusion below $5\%$.}
    \label{fig:sampling-recovery}
\end{figure*}


\subsection{Identifiability}

\begin{theorem}[Identifiability under coherence]
\label{thm:ident}
If the team is in the coherent or phase-locked regime ($\Rens \geq 0.80$)
during the goal build-up, the morphism factorisation is identifiable
from the data $(\Ocal_{\mathrm{ball}}, \Acal)$ up to category
automorphisms.
\end{theorem}

\begin{proof}
By Proposition~\ref{prop:backward}, the morphism factorisation at
phase-lock is unique modulo automorphism. The coherent regime is the
small perturbation of phase-lock; for $\Rens \geq 0.80$, the
Theorem~\ref{thm:template} offset bound is sufficiently tight that the
factorisation remains unique under standard tolerance.
\end{proof}

The theorem motivates pre-filtering build-ups by team coherence regime:
goals scored from incoherent or aperture-dominated regimes are not
amenable to backward template reconstruction because the team did not
have a coherent template to begin with. Such goals must be analysed at
the per-player level rather than at the theoretical-player level.

% ═══════════════════════════════════════════════════════════════════
\section{The World Cup as Reference Dataset}
\label{sec:wc}
% ═══════════════════════════════════════════════════════════════════

\subsection{Corpus properties}

The FIFA World Cup is convenient as a reference corpus for the framework
because of structural properties: (i) 64 matches under uniform rules,
broadcast at uniform technical specifications; (ii) 32 distinct national
teams, each with a recognisable theoretical player developed over
years of preparation; (iii) high cadence of competitive matches over a
short window, supplying longitudinal cell-registry data; (iv)
universally-installed perimeter advertising panels supporting the
homography substrate.

\subsection{Computable quantities}

Within the framework, the World Cup supplies values for the following
quantities per team and per match:
\begin{enumerate}[leftmargin=1.4em]
  \item Coherence trajectory $\Rens(t)$ across each match, with regime
    classification (Definition~\ref{def:regimes}).
  \item Critical coupling $\Kc$ estimate from the empirical
    distribution of player natural gait frequencies.
  \item Per-team tactical generator set (Definition~\ref{def:gen})
    extracted from the goal ledger.
  \item Per-team theoretical-player registry (Definition~\ref{def:theoretical}).
  \item Pairwise team category isomorphism scores, identifying tactical
    families.
\end{enumerate}

\subsection{Single-subject extension}

The framework supports single-subject longitudinal study: a particular
player tracked across multiple matches produces a per-player cell
registry that can be compared to the team registry. The per-player
deviation from the team registry quantifies that player's individuality
within the team template. Star players are expected to exhibit larger
individual deviations than role players.

\subsection{Cross-team category isomorphism}

\begin{definition}[Tactical family]
\label{def:family}
Two teams belong to the same \emph{tactical family} if their role
categories $\Mcal_1, \Mcal_2$ are isomorphic as small categories and
their cell registries agree on the isomorphism to within a tolerance
$\epsilon$.
\end{definition}

The expected outcome of running the framework over an entire World Cup
is the empirical partition of the 32 teams into a small number of
tactical families, with the framework supplying the precise category
isomorphism witnessing the family membership.

% ═══════════════════════════════════════════════════════════════════
\section{Discussion}
% ═══════════════════════════════════════════════════════════════════

\subsection{Scope}

The framework is an \emph{accounting instrument}: it tallies and
structures observed biomechanical activity in a way that respects the
true causal coupling of the team and the inter-body contact. It is not
a predictive engine for match outcomes, a tactics recommender, or an
identification system for player evaluation in the conventional sense.

\subsection{Comparison with existing systems}

Commercial broadcast tracking systems such as TRACAB, STATS
SportVU, and Hawk-Eye football tracking estimate per-player
$(x,y)$ pitch positions at 25--50 Hz \citep{theagarajan2021,dorazio2009}.
These systems supply the formation matrix $\Phi(t)$ used here but do
not provide foot-contact biomechanics, do not segment duels as cognate
pairs, do not compute coherence regimes, and do not perform morphism
reconstruction. The proposed framework is complementary to existing
tracking: it consumes the formation matrix where available and
substitutes the panel-homography channel where it is not.

\subsection{Limitations}

Three limitations bound the framework's near-term applicability.

\textbf{Computer vision pipeline.} Foot-contact detection from
broadcast footage with the panel-homography substrate is novel; while
the substrate's geometric basis is sound, robust implementation in
the presence of motion blur, occlusion, and broadcast cut transitions
is non-trivial and remains to be demonstrated at scale.

\textbf{Identity disambiguation.} The transcript-driven identity
assignment is robust to the typical 1--3 s commentary lag, but
sustained micro-passes (one-touch interplay) below the commentary
resolution may not be individually attributable. Aggregation over
several touches typically resolves this.

\textbf{Goalkeeping.} Goalkeepers occupy a degenerate role: they are
ball-channel high-fidelity samples but participate in the field
coupling only sporadically. Whether to include them in the
$N{=}11$ Kuramoto ensemble or to model them as an exogenous boundary
condition is a modelling choice we leave open.

\subsection{Relation to tactical theory}

The framework is consistent with long-standing observations from match
analysis \citep{reep1968,hughes2005,memmert2017} that elite football is
finitely generated at the tactical level. It supplies a formal mechanism
for what analysts have called ``team identity'' or ``playing
philosophy'' and quantifies the regime of team coherence under which
that identity is operationally valid.

% ═══════════════════════════════════════════════════════════════════
\section{Conclusion}
% ═══════════════════════════════════════════════════════════════════

We have proposed an observation-operator framework for the biomechanical
accounting of association football. The team is treated as a partially
observable bioreactor of eleven coupled biomechanical cells; the ball is
the primary high-fidelity sample probe; the broadcast camera supplies
the low-resolution ambient indicator. The eleven cells couple via
Kuramoto-type phase dynamics, and the team's coherence is measured by
the order parameter $\Rens$. The five regimes of $\Rens$ partition match
play into operationally distinct collective modes, and the phase-locked
regime is precisely the regime in which a single theoretical player ---
formalised as the team's role category together with a cell registry ---
correctly represents the collective.

Duels are formalised as cognate oscillator pairs, the irreducible
accounting unit at which contact-coupled mechanical load is correctly
totalled. A measurement substrate based on lower-body tracking against
the homography of stadium advertising panels delivers metric pitch
coordinates and player/spectator segmentation in a single geometric
step, side-stepping the resolution failures of full-skeleton pose
estimation on broadcast footage. Goal sequences are reconstructed
backwards as morphism compositions in the role category, recovering
both the morphism factorisation and the cell registry entries that
realise it.

The FIFA World Cup, with 64 matches and 32 distinct national templates
across uniform broadcast conditions, supplies a near-ideal reference
corpus. The expected output is a structured biomechanical ledger per
team---coherence trajectories, generator sets, theoretical-player
registries, and tactical-family memberships---not a predictive model.
The framework supplies a precise vocabulary for what match analysts
have long described informally: team identity, cohesion, and the
recurrent generative patterns that produce goals.

\section*{Data and code availability}

No data are reported in the present article; the framework is
theoretical. Reference implementations of the homography substrate and
the morphism-reconstruction algorithm are available from the author
on request.

\bibliographystyle{unsrtnat}
\bibliography{references}

\end{document}
